\section{Many constraints: the larger case and its classical preprocessing}\label{sec:large}
This appendix gives the multi-constraint case of \cref{sec:reduction} in full: the case, the width of the coherent state, the classical preprocessing derived from the automaton, quality against penalty QAOA, circuits for a device, the block scheme for $100$--$200$ variables and the fault-tolerant estimate.

\paragraph{The case.} A multi-period long/short portfolio with $A$ assets (grouped in sectors) and $T$ periods; the decision variables are $x_{t,i,d}$ with $d\in\{\mathrm{long},\mathrm{short}\}$. Per period: capital $0\le L_t-S_t\le3$ and budget $L_t+S_t\le3$ (as in \cref{sec:portfolio}), at most one long position per sector, and no long and short position on the same asset; globally, a trading budget $\sum_t(L_t+S_t)\le Q$ with $Q=\lceil 0.6\cdot3T\rceil$. With $T=8$, $A=12$ and three sectors this is $n=192$ variables and $137$ constraints. The automaton is written from the structure of the problem (\path{negsearch/large_case.py}) and validated against brute-force enumeration (at $n=16$ both variable orders give exactly the $521$ feasible strings); the same compiler as before then applies. The objective is quadratic with couplings across assets and periods (risk and a transaction term), so a dynamic programme over the automaton does not solve it.

\paragraph{The coherent state does not fit hardware.} \Cref{tab:large_scaling} shows that the coherent state $A_q$ over all constraints has width $11$ at $n=24$ and $103$ at $n=192$ (the global budget couples all periods), and its compiled size grows from $12.9$ thousand to $5.1$ million CZ gates and from $88$ to $1156$ qubits (default synthesis, before routing). The uniform state has a feasible share of $2^{-12.6}$ at $n=24$ and $2^{-131}$ at $n=192$, which is what the construction fixes; but it uses \emph{more} qubits than a slack-variable QUBO ($252$ at $n=192$).

\begin{table*}[t]
\centering\footnotesize\setlength{\tabcolsep}{2.6pt}
\caption{\orig{exact automaton construction and compiled counts} Multi-period portfolio with sector caps, long/short exclusivity and a global trading budget ($A$ assets, $T$ periods, sectors of $3$--$5$ assets). Left: the coherent state $A_q$ over all constraints (binary register, period-major order). Right: classical preprocessing that conditions on the interface. CZ counts (thousands) use Qiskit's default ancilla-free multi-controlled synthesis, before routing (routing roughly doubles them). ``Per-period'' conditions on the number of positions of each period (binary registers, sum over periods); ``per-sector'' also on the counts of each sector (one-hot registers; typical total over the blocks, and the largest single block in CZ). The uniform column is the share of feasible strings in the uniform state; the slack column is the number of qubits of a slack-variable QUBO.}\label{tab:large_scaling}
\begin{tabular}{lrrrrrrrrrr}
\toprule
 & & & \multicolumn{3}{c}{coherent $A_q$} & & & \multicolumn{3}{c}{conditioned (k CZ)}\\
\cmidrule(lr){4-6}\cmidrule(lr){9-11}
$A,T$ & $n$ & constr. & width & qubits & kCZ & unif. & slack q. & per-period & per-sector & max block\\
\midrule
6,2 & 24 & 21 & 11 & 88 & 12.9 & $2^{-13}$ & 39 & 1.3 & 0.3 & 115\\
8,3 & 48 & 37 & 27 & 207 & 116.5 & $2^{-30}$ & 69 & 3.4 & 0.7 & 209\\
10,4 & 80 & 57 & 31 & 380 & 339.4 & $2^{-53}$ & 107 & 6.3 & 1.7 & 349\\
10,6 & 120 & 85 & 52 & 636 & 1,189.6 & $2^{-77}$ & 160 & 9.4 & 2.5 & 349\\
12,8 & 192 & 137 & 103 & 1156 & 5,058.5 & $2^{-131}$ & 252 & 29.0 & 2.9 & 209\\
\bottomrule
\end{tabular}
\end{table*}


\paragraph{Classical preprocessing from the theory.} Three steps, all computed before any circuit is run. (i)~\emph{Variable order.} The width of the exact automaton depends strongly on the order: period-major gives $8$, $23$ and $11$ states against $36$, $172$ and $36$ for asset-major on three small instances. (ii)~\emph{Conditioning on an interface.} If the number of positions of each period is drawn classically, the periods become independent blocks and the cost grows linearly in $T$: from $12.9$ to $1.3$ thousand CZ at $n=24$, and from $5.1$ million to $29$ thousand at $n=192$. (iii)~\emph{A second level.} Conditioning also on the counts per sector, with one-hot registers (which beat binary registers on blocks of width $\ge5$ by a factor $3$--$5$), leaves blocks of at most $349$ CZ each and a typical total of $2.9$ thousand CZ at $n=192$. This is a decomposition by separators, as in tree decompositions: the generalisation of ``draw the counts classically and prepare Dicke states'' of \cref{sec:portfolio}. The initial distribution of the interface must have full support on the feasible interface values (we use weights proportional to the number of completions).

\paragraph{More preprocessing before the circuit.} Three further passes act on the constraint set before the automaton is compiled (\path{negsearch/preprocess.py}, \path{negsearch/symlower.py}; \cref{tab:large_order,tab:large_symmetry}). The first two leave the feasible set unchanged, up to the forced variables, so feasibility by construction is preserved. (i)~\emph{Cleaning and order search.} Repeated and subsumed clauses are removed, unit clauses are propagated (a forced variable stops being a qubit), and the variable order is chosen to minimise the automaton width by hill climbing on adjacent swaps. On random conflict constraints with $14$--$26$ variables the width falls from $10$, $20$, $20$ and $269$ (index order) to $5$, $14$, $18$ and $49$, and a low-degree-first order is far worse ($24$ to $188$); the clause counts in \cref{tab:large_order} come from injected redundancy and do not estimate natural instances. (ii)~\emph{Symmetry.} In every block of the portfolio the (long, short) pairs of the assets are interchangeable, which an automatic check on the feasible set detects. The block is then a set of arrangements, and instead of the generic register-based $A_q$ it is prepared with Dicke states on the long qubits and on the compressed short qubits, followed by a chain of controlled swaps that inserts the long positions. The state is exactly uniform on the feasible set (checked for every block with up to five assets, $l\le1$ and $s\le3$), it imposes the exclusivity, and it needs no register, no uncomputation and no post-selection. It cuts the compiled size by $3.1$--$4.4\times$ against the generic circuit (for five assets with $(l,s)=(1,1)$: $598\to160$ CZ and $43\to10$ qubits) and by $5$--$8\times$ when the generic registers must be uncomputed, which QAOA with pair-exchange mixers requires. For plain counts a Dicke state alone is cheaper still ($16$ CZ for the same block, without the exclusivity): on symmetric counting structure the generic automaton is not the competitive circuit, and its role is to find the structure and to give the support guarantee. (iii)~\emph{Separators}, as above.

\begin{table*}[t]
\centering\footnotesize\setlength{\tabcolsep}{4pt}
\caption{\orig{exact automaton construction} Variable order and automaton width for conflict constraints (random pair-exclusion graphs with repeated, subsumed and unit clauses injected). Width of the reduced automaton for the index order, a low-degree-first order, a degeneracy order, a colouring-layer order and the best order found by hill climbing on adjacent swaps (starting from the best of the four). Cleaning removes repeated and subsumed clauses and propagates unit clauses (forced variables stop being qubits); the repeated and subsumed clauses are injected, so the reduction in clauses is not an estimate for natural instances.}\label{tab:large_order}
\begin{tabular}{rrrrrrrr}
\toprule
$n$ & clauses (raw$\to$clean) & qubits & index & low degree & degeneracy & colouring & searched\\
\midrule
14 & 38$\to$20 & 12 & 10 & 24 & 12 & 11 & 5\\
18 & 52$\to$29 & 17 & 20 & 66 & 23 & 23 & 14\\
22 & 42$\to$23 & 19 & 20 & 120 & 20 & 56 & 18\\
26 & 64$\to$39 & 25 & 269 & 188 & 100 & 182 & 49\\
\bottomrule
\end{tabular}
\end{table*}

\begin{table*}[t]
\centering\footnotesize\setlength{\tabcolsep}{3.5pt}
\caption{\orig{compiled counts} Symmetry-aware lowering of $A_q$ for a block of $m$ interchangeable assets conditioned on $l$ longs and $s$ shorts (all units found interchangeable by the automatic check). CZ after routing to \texttt{ibm\_fez} (best of three seeds). The generic $A_q$ uses a one-hot automaton register, relative-phase Toffolis, and for QAOA the registers uncomputed; the lowering has no register, no uncomputation and no post-selection.}\label{tab:large_symmetry}
\begin{tabular}{rrrrrrrr}
\toprule
$m$ & $(l,s)$ & $|F|$ & generic: qubits & CZ & CZ uncomputed & lowering: qubits & CZ\\
\midrule
3 & (1,1) & 6 & 21 & 190 & 281 & 6 & 46\\
4 & (1,1) & 12 & 32 & 385 & 733 & 8 & 95\\
4 & (1,2) & 12 & 33 & 430 & 773 & 8 & 112\\
5 & (0,3) & 10 & 30 & 269 & 457 & 10 & 86\\
5 & (1,1) & 20 & 43 & 598 & 1136 & 10 & 160\\
5 & (1,2) & 30 & 49 & 872 & 1590 & 10 & 200\\
\bottomrule
\end{tabular}
\end{table*}


\paragraph{What the preprocessing costs.} Drawing the interface classically replaces a superposition over interface values by a mixture. This is valid for sampling and for QAOA with mixers that conserve the interface: the Grover mixer applied to each block separately keeps every block's constraints, hence every sample feasible (case~(b) of \cref{thm:qaoa}). It is \emph{not} valid for amplitude amplification of the whole state or for the global Grover mixer, which need coherence across interface values. It also moves the quantum content: the initial distribution of a conditioned state is a product of tiny blocks that a classical computer samples trivially, and the quantum part is the cost layer, whose quadratic couplings entangle the blocks, and the block mixers.

\paragraph{Quality against penalty and unbalanced QAOA.} At $n=16$ (\cref{tab:large_quality}) the state keeps every sample feasible, and the conditioned version, with the interface drawn classically and a Grover mixer per interface value, loses very little against the coherent one (quality $0.71$ against $0.73$ at $p=3$). Penalty QAOA and unbalanced QAOA, with a uniform start, an $X$ mixer and the better of two penalty weights, reach $0.17$ and $0.44$ at $p=3$, with $39$\,\% and $83$\,\% feasible samples. The study is small: two random objectives, locally optimised angles from three restarts, and $P(\mathrm{opt})$ at most $0.02$ for every method because $521$ strings are feasible, so the gap is in quality and feasibility, not in finding the optimum.

\begin{table}[t]
\centering\small\setlength{\tabcolsep}{4pt}
\caption{\orig{exact ideal numerics} Quality on the multi-constraint case at $n=16$ (4 assets, 2 periods, 2 sectors; $521$ feasible strings; uniform feasible share $0.8$\,\%; mean over 2 random objectives). QAOA depth $p$, exact numerics with locally optimised angles. Quality counts infeasible samples as $0$; $P(\mathrm{opt})$ is the probability of an optimal string (uniform feasible sampling: 0.0019). Preparation alone ($p=0$): feasible $1.00$, quality 0.44.}\label{tab:large_quality}
\begin{tabular}{llrrr}
\toprule
$p$ & method & feasible & quality & $P(\mathrm{opt})$\\
\midrule
1 & $A_q$ + Grover mixer & 1.00 & 0.61 & 0.008\\
1 & conditioned + Grover mixer & 1.00 & 0.61 & 0.008\\
1 & penalty QAOA & 0.34 & 0.16 & 0.000\\
1 & unbalanced QAOA & 0.65 & 0.26 & 0.000\\
\midrule
2 & $A_q$ + Grover mixer & 1.00 & 0.69 & 0.014\\
2 & conditioned + Grover mixer & 1.00 & 0.69 & 0.012\\
2 & penalty QAOA & 0.24 & 0.15 & 0.003\\
2 & unbalanced QAOA & 0.73 & 0.41 & 0.005\\
\midrule
3 & $A_q$ + Grover mixer & 1.00 & 0.73 & 0.019\\
3 & conditioned + Grover mixer & 1.00 & 0.71 & 0.016\\
3 & penalty QAOA & 0.39 & 0.17 & 0.001\\
3 & unbalanced QAOA & 0.83 & 0.44 & 0.001\\
\bottomrule
\end{tabular}
\end{table}


\paragraph{Circuits for hardware.} We generated and validated circuits for two instances (\path{experiments/hardware_large_case.py}; \cref{tab:large_hw}). For $n=24$ (six assets, two periods, two sectors; $21$ constraints, $2620$ feasible strings among $1.7\cdot10^7$) the preparation of one interface tuple uses $150$--$184$ CZ, depth $136$--$238$ and $40$--$42$ qubits, with an estimated success probability of $0.39$--$0.46$ and a duration of about $0.1\,T_2$. The QAOA layer is not executable at this size: the cost layer (about $640$ CZ after routing) and four block mixers bring it to $3.8$--$4.4$ thousand CZ. At $n=8$ (four assets, one period) a full $p=1$ QAOA circuit has $443$--$488$ CZ, depth $693$--$775$, $16$ qubits and an estimated success probability of $0.18$--$0.2$. Validation: the logical $n=8$ QAOA circuit reproduces the exact product-space numerics to $6\cdot10^{-16}$ for all eight interface tuples, and the $n=24$ preparation, simulated as matrix product states, gives $100\,\%$ feasible samples and a distance to the exact distribution within shot noise (total variation $0.023$--$0.025$ for supports of $81$ strings, $20{,}000$ shots). At larger $n$ the dense quadratic layer and the block mixers set the limit, not the preparation. The script submits the circuits to a device (with a minimum estimated success probability), and \path{experiments/hardware_large_case_report.py} reports the share of feasible samples with Wilson intervals; no data from a device are reported here.

\begin{table*}[t]
\centering\small\setlength{\tabcolsep}{4pt}
\caption{\orig{compiled counts and ESP (model); nothing executed} Hardware circuits generated for the conditioned multi-constraint case, compiled on the \texttt{ibm\_fez} topology (FakeFez snapshot; best of 10 transpiler seeds by estimated success probability; ranges over the interface tuples compiled). Nothing was executed on a device. The QAOA layer is not executable at $n=24$ (about $3.8$--$4.4$ thousand CZ, estimated success probability $\approx0$), so only the preparation was generated there.}\label{tab:large_hw}
\begin{tabular}{llrrrrr}
\toprule
instance & circuit & qubits & CZ & depth & $\mu$s & ESP\\
\midrule
$n=24$ & $A_q$ only & 40--42 & 150--184 & 136--238 & 6.4--10.3 & 0.39--0.46\\
$n=8$ & $A_q$ only & 12 & 20--22 & 34 & 2.7 & 0.89--0.90\\
$n=8$ & $A_q$ + QAOA $p{=}1$ & 16 & 443--488 & 693--775 & 28.1--31.5 & 0.18--0.20\\
\bottomrule
\end{tabular}
\end{table*}


\paragraph{Ideal, optimised and executable (QAOA).} \Cref{tab:large_qaoa_real} compares the first version of the QAOA layer with the optimised one, compiled to the \texttt{ibm\_fez} topology. Pair-exchange mixers instead of a Grover mixer per block cut the CZ count by $1.8$--$2.1\times$ ($443\to243$ at $n=8$, $1391\to660$ at $n=16$, $4002\to1894$ at $n=24$): in the first version the reflection of the Grover mixer carried about $60\%$ of the gates, and the dense objective layer ($28$ couplings at $n=8$, $156$ at $n=24$) is intrinsic to the quadratic objective. Pruning the objective couplings below $25\%$ of the largest cuts a further $17$--$43\%$ and leaves feasibility untouched; its cost is a small loss of quality ($0.574\to0.567$ at $n=8$ and $0.541\to0.531$ at $n=16$, against $0.514$ and $0.474$ for the uniform feasible preparation) and of the probability of the optimum ($0.0046\to0.0015$ at $n=16$, $p=1$). The result: one layer at $n=8$ is estimated to be within reach of current devices (success probability $0.45$--$0.54$) and improves on the preparation alone by $0.06$ in quality; $n=12$ and $n=16$ are on the frontier ($0.12$ and $0.29$, a weak signal); and at $n=24$ no variant is executable ($0.01$ with all optimisations). The quality at $n=16$ is exact and averaged over all interface tuples, with angles from four restarts. A noisy device returns a mixture of this distribution and noise, so the usable gain shrinks with the estimated success probability.

\begin{table*}[t]
\centering\footnotesize\setlength{\tabcolsep}{4pt}
\caption{\orig{compiled counts, ESP (model) and exact ideal simulation; nothing executed} Multi-constraint case, one QAOA layer ($p=1$), one interface tuple, compiled to the \texttt{ibm\_fez} topology (best of eight transpiler seeds; nothing executed). ESP is the product of one minus the error of every gate. Status: ``run today'' ESP $\ge0.3$, ``frontier'' $0.05$--$0.3$, ``out of reach'' below $0.05$. Quality is measured on the exact feasible set, averaged over all interface tuples (higher is better, $1$ is the optimum); the preparation alone is the uniform feasible state. Pruning drops the objective couplings below $25\%$ of the largest, which never affects feasibility. The pair-exchange mixer $e^{-i\beta P}$ with $P=\mathrm{SWAP}(\mathrm{long}_j,\mathrm{long}_k)\,\mathrm{SWAP}(\mathrm{short}_j,\mathrm{short}_k)$ conserves all counts and the exclusivity and needs the preparation registers uncomputed; the preparation-only row is that circuit without the QAOA layer (generic $A_q$, registers uncomputed).}\label{tab:large_qaoa_real}
\begin{tabular}{llrrrrrll}
\toprule
$n$ & variant & qubits & CZ & depth & ESP & $d/T_2$ & quality & status\\
\midrule
8 & Grover mixer (first version) & 16 & 443 & 774 & 0.21 & 0.31 & 0.574 & frontier\\
 & pair-exchange mixers & 12 & 243 & 385 & 0.45 & 0.13 & 0.574 & run today\\
 & \quad + objective pruning 25\% & 12 & 190 & 310 & 0.54 & 0.10 & 0.567 & run today\\
 & preparation only & 10 & 44 & 71 & 0.88 & 0.02 & 0.514 & run today\\
\midrule
12 & Grover mixer (first version) & 28 & 1\,481 & 2\,183 & $2.0\!\cdot\!10^{-3}$ & 1.36 & -- & out of reach\\
 & pair-exchange mixers & 22 & 696 & 1\,037 & 0.08 & 0.38 & -- & frontier\\
 & \quad + objective pruning 25\% & 20 & 579 & 881 & 0.12 & 0.33 & -- & frontier\\
 & preparation only & 20 & 158 & 293 & 0.61 & 0.10 & -- & run today\\
\midrule
16 & Grover mixer (first version) & 29 & 1\,391 & 1\,551 & $2.9\!\cdot\!10^{-3}$ & 0.94 & 0.541 & out of reach\\
 & pair-exchange mixers & 24 & 660 & 599 & 0.07 & 0.35 & 0.541 & frontier\\
 & \quad + objective pruning 25\% & 24 & 378 & 318 & 0.29 & 0.10 & 0.531 & frontier\\
 & preparation only & 20 & 88 & 73 & 0.75 & 0.02 & 0.474 & run today\\
\midrule
24 & Grover mixer (first version) & 48 & 4\,002 & 4\,143 & $4.3\!\cdot\!10^{-153}$ & 2.10 & -- & out of reach\\
 & pair-exchange mixers & 40 & 1\,894 & 1\,202 & $9.6\!\cdot\!10^{-5}$ & 0.58 & -- & out of reach\\
 & \quad + objective pruning 25\% & 43 & 1\,202 & 1\,144 & $1.0\!\cdot\!10^{-2}$ & 0.47 & -- & out of reach\\
 & preparation only & 38 & 343 & 318 & 0.29 & 0.12 & -- & frontier\\
\bottomrule
\end{tabular}
\end{table*}


\paragraph{Problems of 100--200 variables as blocks.} The way to put a problem of this size on a device is many circuits, not one (\cref{tab:large_blocks}): with the interface drawn classically, each (period, sector) block is a $p=1$ circuit of $8$--$10$ qubits with the other blocks fixed as linear fields, and a classical loop sweeps the blocks. At $n=120$ and $192$ every block has an estimated success probability of at least $0.20$ (median $0.32$ and $0.47$) and $144$--$343$ CZ gates at the median, so a sweep is $12$ or $24$ circuits and $2.9$--$5.8$ thousand CZ in total, against $1.2$ and $5.1$ million CZ for the coherent state of \cref{tab:large_scaling}. In exact simulation with ten shots per block, per-block QAOA closes $96$--$102\%$ of the gap between a random feasible assignment and the block optimum, but uniform sampling closes $94$--$99\%$: a block has at most a few hundred feasible strings and is solved by enumeration, so this shows that the circuits are executable, not that QAOA helps inside them.
\begin{table*}[t]
\centering\footnotesize\setlength{\tabcolsep}{3.5pt}
\caption{\orig{compiled counts and ESP (model); nothing executed} QAOA as many small blocks. The interface is drawn classically and every (period, sector) block is a separate $p=1$ circuit (pair-exchange mixers, preparation from the symmetry-aware lowering of $A_q$, exact and without registers), while the other blocks are held fixed and enter as linear fields. Compiled to the \texttt{ibm\_fez} topology, three random interface tuples per instance; CZ and ESP are over all blocks.}\label{tab:large_blocks}
\begin{tabular}{lrrlrrrr}
\toprule
$A,T$ & $n$ & blocks & objective & qubits/block & CZ/block (med.\,/\,max) & ESP (med.\,/\,min) & CZ per sweep\\
\midrule
10,6 & 120 & 12 & full & 10 & 343\,/\,505 & 0.32\,/\,0.20 & 4\,471\\
10,6 & 120 & 12 & pruned 25\% & 10 & 205\,/\,394 & 0.49\,/\,0.26 & 2\,901\\
12,8 & 192 & 24 & full & 8 & 235\,/\,324 & 0.47\,/\,0.34 & 5\,766\\
12,8 & 192 & 24 & pruned 25\% & 8 & 144\,/\,248 & 0.62\,/\,0.46 & 3\,583\\
\bottomrule
\end{tabular}
\end{table*}


\paragraph{Amplitude amplification: a fault-tolerant estimate.} The setting where amplification over the coherent state applies is the monolithic one, and \cref{tab:large_grover_ft} gives its cost on a fault-tolerant machine under a simple model: $88$ to $1156$ logical qubits and $3.9\cdot10^3$ to $6.8\cdot10^5$ Toffoli gates per iteration. Amplification beats classical sampling from the same state only when the target set is a tiny fraction $a$ of the feasible set, below $a^\star$ ($2\cdot10^{-7}$ to $2\cdot10^{-11}$ at $n=24$, $7\cdot10^{-12}$ to $7\cdot10^{-16}$ at $n=192$). For $a=10^{-12}$ and $n=24$ the model gives $6$ hours to $2.5$ days against $12$ to $116$ days of classical sampling, while at $n=192$ it gives $43$ to $435$ days and no advantage. The value of $a$ is a free parameter here, not computed for this instance, and the model leaves out the threshold oracle, the overhead of error correction and magic-state production; it is a statement about the machine this use would need, and it is not a measurement.

\begin{table*}[t]
\centering\footnotesize\setlength{\tabcolsep}{3pt}
\caption{\orig{resource estimate for a fault-tolerant machine; not a measurement} Amplitude amplification over the coherent state $A_q$ on a fault-tolerant machine (a model, not a measurement). A multi-controlled gate with $c$ controls costs $2(c-1)$ Toffoli gates, a Toffoli $7$ T gates, and one iteration is $A_q^\dagger$, the reflection and $A_q$ (the threshold oracle is common to all methods and not counted). With success probability $a$ the amplification needs $\frac\pi4/\sqrt a$ iterations and classical sampling from the same state needs $1/a$ samples. Times assume $1$--$10\,\mu$s per T gate and $1$--$10\,\mu$s per classical sample. $a^\star$ is the success probability below which amplification is faster than classical sampling.}\label{tab:large_grover_ft}
\begin{tabular}{lrrrrr}
\toprule
$A,T$ & logical qubits & Toffoli / iteration & $a^\star$ (optimistic\,/\,conservative) & time at $a=10^{-9}$ & time at $a=10^{-12}$\\
\midrule
6,2 & 88 & 3\,954 & $2.1\!\cdot\!10^{-7}$\,/\,$2.1\!\cdot\!10^{-11}$ & 11\,min--1.9\,h & 6.0\,h--2.5\,d\\
8,3 & 207 & 23\,736 & $5.9\!\cdot\!10^{-9}$\,/\,$5.9\!\cdot\!10^{-13}$ & 1.1\,h--11.5\,h & 36.2\,h--15\,d\\
10,4 & 380 & 63\,710 & $8.2\!\cdot\!10^{-10}$\,/\,$8.2\!\cdot\!10^{-14}$ & 3.1\,h--30.8\,h & 4.1\,d--41\,d\\
12,8 & 1\,156 & 682\,838 & $7.1\!\cdot\!10^{-12}$\,/\,$7.1\!\cdot\!10^{-16}$ & 33.0\,h--14\,d & 43\,d--435\,d\\
\bottomrule
\end{tabular}
\end{table*}


\paragraph{What this establishes.} The construction applies to a case with $137$ constraints of four kinds. Classical preprocessing from the automaton (order, interface conditioning, symmetry) reduces the circuit by orders of magnitude, from $5.1$ million to a few thousand CZ at $n=192$, at the price of coherence across the interface. By compiled counts and the noise model the picture is graded: the preparation of a conditioned state is predicted to be executable up to $n=24$ ($150$--$184$ CZ with the generic compiler, ESP $0.39$--$0.46$) and one QAOA layer up to $n=8$, with $n=12$--$16$ on the frontier; problems of $120$--$192$ variables are reachable only as blocks of at most ten qubits, where QAOA does not beat sampling; and amplitude amplification over the coherent state is a fault-tolerant use, with its own cost model. We do not claim an advantage over classical solvers; a mixed-integer solver handles these sizes, and the comparisons of quality are against other quantum encodings.
